% =====================================================================
%  methods-v2 / proposed-solution-blue.tex
%  Manuscript-ready DRAFT of the Proposed Solution section.
%  Deliverable for MICROENCODER_KD_BRIEF.md Sec. 20.
%  Prepared 2026-10-02 (gnhf microencoder-kd sprint, iteration 19).
%
%  WHAT THIS IS
%  ------------
%  Drop-in LaTeX for the empty \section{Proposed Solution} at line 193 of
%  source/main-original.tex. It specifies the distilled micro-encoder (TLM)
%  router: f_theta, the KD training, the runtime cost integration, the
%  inference routing rule, and Split-LLM integration. Content is lifted
%  verbatim from the locked methods-v2 deliverable 06 (proposed-method),
%  which in turn rests on 00-05, 07-09.
%
%  BLUE CONVENTION (brief Sec. 20)
%  -------------------------------
%  All newly proposed manuscript prose is wrapped in \prop{...}, which
%  renders in blue. Shared/background sentences that merely restate the
%  already-accepted problem formulation are left in black. Equations and
%  algorithms that are new are flagged in blue in their lead-in sentence.
%
%  HOW TO MERGE INTO THE IEEEtran MANUSCRIPT (do NOT edit the original here)
%  ----------------------------------------------------------------------
%    1. Add to the manuscript preamble (it is not currently loaded):
%           \usepackage{xcolor}
%           \newcommand{\prop}[1]{\textcolor{blue}{#1}}
%    2. Replace the stub body of \section{Proposed Solution} (the lone ".")
%       with the body between the BEGIN/END MERGE markers below.
%    3. The equations below are written with manuscript notation
%       (\mathcal{M}, M_k, f_\theta, \lambda, C) so they slot in after the
%       existing Eqs. (1)-(3) in Sec. "Problem Statement and Formulation".
%
%  COMPILATION STATUS
%  ------------------
%  This file is a STANDALONE, self-contained document (article class) so it
%  can be compiled on its own to proof the prose/equations/algorithms. No
%  LaTeX toolchain exists in the sprint environment, so it has NOT been
%  compile-verified here (same honest caveat as architecture-figure.tex,
%  iter 16); it was authored to compile but must be checked at merge time.
%
%  EVIDENCE DISCIPLINE (brief Sec. 23)
%  -----------------------------------
%  No results are claimed. Per brief Sec. 20/26 this file contains no
%  "experiments demonstrate / we outperform / our method achieves"
%  language; every performance statement is phrased as a design intent or
%  a to-be-evaluated hypothesis. Bracketed [S:NN] comments trace each claim
%  to its source deliverable; they are LaTeX comments and never render.
% =====================================================================
\documentclass[11pt]{article}
\usepackage[margin=1in]{geometry}
\usepackage{amsmath,amsfonts,amssymb}
\usepackage{algorithm}
\usepackage{algorithmic}
\usepackage{xcolor}
\usepackage{hyperref}

% ---- blue convention for newly proposed manuscript prose (brief Sec. 20) ----
\newcommand{\prop}[1]{\textcolor{blue}{#1}}

% ---- notation shorthands (match source/main-original.tex) ----
\newcommand{\Mcal}{\mathcal{M}}
\newcommand{\Qcal}{\mathcal{Q}}
\newcommand{\Acal}{\mathcal{A}}
\newcommand{\Uhat}{\widehat{U}}
\newcommand{\argmax}{\operatorname*{arg\,max}}

\title{Proposed Solution (BLUE DRAFT)\\
\large Distilled Micro-Encoder Router for Edge SLM / Split-LLM Collaborative Inference}
\author{methods-v2 draft --- not for submission as-is}
\date{2026-10-02}

\begin{document}
\maketitle

\noindent\textbf{Reviewer note (remove at merge):} blue text is newly
proposed manuscript prose; black text restates the already-accepted problem
formulation from Sec.~``Problem Statement and Formulation''. No experimental
results are stated anywhere in this draft.

% =====================================================================
% ========================  BEGIN MERGE BLOCK  ========================
% =====================================================================

\section{Proposed Solution}
\label{sec:proposed}

\prop{%
We implement the routing function $f_{\theta}$ of
Eq.~(3) as a \emph{distilled micro-encoder}: a Tiny Language Model (TLM)
that predicts, for each candidate execution path, a calibrated utility for
the query, and is trained once with knowledge distillation from a larger
same-family encoder. At inference the TLM runs a single forward pass on the
edge device and its learned utilities are combined with \emph{measured} live
system cost to select an execution path. The design keeps the learned,
query-dependent quality signal separate from the measured, network-dependent
cost signal: the former is trained into the TLM weights, the latter is read
from the running system and never enters the weights.}

% ------------------------------------------------------------------
\subsection{Execution paths and the routing decision}
\label{subsec:actions}

\prop{%
The decision in Eq.~(1) is taken over execution \emph{actions} rather than
over models alone, because a Split-LLM path is defined by a cut depth and two
cut depths of the same model are two different cost/quality operating points.
We therefore route over the action set}
\begin{equation}
\label{eq:actionset}
\Acal \;=\; \bigl\{\, \textsc{Edge},\;
\textsc{Split}(d_1),\dots,\textsc{Split}(d_m),\;
\textsc{Strong}\,\bigr\},
\end{equation}
\prop{where \textsc{Edge} runs the small model entirely on the edge device;
$\textsc{Split}(d)$ runs layers $[0,d)$ on the edge, transfers the
intermediate activation to the server, and runs layers $[d,L)$ there; and
\textsc{Strong} executes the strong/remote model in full. We write
$N=\lvert\Acal\rvert$ for the number of actions, fixed by the model pool and
the enabled split depths.}

\prop{%
Equation~(\ref{eq:actionset}) is an \emph{extension of notation, not a
redefinition of the problem}. Each action $a\in\Acal$ induces a
model-of-record $M(a)$; the manuscript's model set $\Mcal$ is the image
$\{M(a):a\in\Acal\}$, and when split depths are disabled ($m=0$) the map
$a\mapsto M(a)$ is a bijection, so $\Acal\equiv\Mcal$. Under that condition,
and with the cost term switched off ($\lambda=0$), the rule introduced below
collapses exactly to Eq.~(3). The proposed method thus refines Eq.~(3) by
(i)~enriching the choice set with split depths and (ii)~factoring the score
into a learned term and a measured term.}

% ------------------------------------------------------------------
\subsection{The micro-encoder scoring function $f_{\theta}$}
\label{subsec:ftheta}

\prop{%
The TLM is an encoder $g_{\theta}:\Qcal\!\rightarrow\!\mathbb{R}^{N}$ with a
pooled query representation feeding $N$ parallel scalar heads, one per
action. For a query $q$,}
\begin{align}
\label{eq:pool}
h(q) &\;=\; \mathrm{Pool}\bigl(\mathrm{Encoder}_{\theta}(q)\bigr)
      \;\in\; \mathbb{R}^{H},\\[2pt]
\label{eq:head}
\Uhat_{\theta}(a\mid q) &\;=\; \sigma\bigl(w_a^{\top} h(q) + b_a\bigr),
      \qquad a\in\Acal,
\end{align}
\prop{%
where $\sigma$ is the logistic function and each head is \emph{independent}
(a per-action sigmoid, not a softmax over actions). Independent heads are a
deliberate choice: they keep the learned utility separable from the cost term
introduced below, and they make the model pool extensible, since adding an
action is adding one head rather than re-training a joint normalisation over
all actions. The proposed scoring function of Eq.~(3) is the per-action
learned utility}
\begin{equation}
\label{eq:ftheta}
f_{\theta}(q,a) \;:=\; \Uhat_{\theta}(a\mid q),
\end{equation}
\prop{%
i.e.\ the preferred $f_{\theta}(q,M_k)=\Uhat_{\theta}(M_k\mid q)$ form, with
the model index $M_k$ generalised to the action $a$ of
Eq.~(\ref{eq:actionset}). The target semantics of $\Uhat_{\theta}(a\mid q)$
is a \emph{calibrated probability that action $a$ succeeds on query $q$}
(Sec.~\ref{subsec:labels}).}

\prop{%
We instantiate $g_{\theta}$ with a compact modern encoder backbone in the
$\sim$20--150\,M parameter range (the micro-encoder), exported to a quantised
(e.g.\ INT8/ONNX) runtime for edge execution; the exact backbone and
precision are an implementation choice to be fixed by ablation rather than a
contribution of this work.}
% [S:05] backbone size-range and INT8/ONNX edge path; exact backbone is [OPEN].

% ------------------------------------------------------------------
\subsection{Runtime routing rule: learned utility minus measured cost}
\label{subsec:rule}

\prop{%
At inference, for query $q$ under measured live system state $s$, the router
selects}
\begin{equation}
\label{eq:rule}
R(q,s) \;=\; \argmax_{a\in\Acal}\;
\Bigl[\, \underbrace{\Uhat_{\theta}(a\mid q)}_{\text{learned (TLM)}}
\;-\; \lambda\cdot \underbrace{C(a\mid s)}_{\text{measured, not in }\theta}
\,\Bigr],
\end{equation}
\prop{%
where $\lambda$ is the quality--cost trade-off coefficient of Eq.~(2).
$\Uhat_{\theta}(a\mid q)$ is produced by a single forward pass of the TLM and
is \emph{system-unaware}: it depends on the query, not on the network.
$C(a\mid s)$ is the scalarised live system cost of action $a$ under the
measured state $s$, where $s$ aggregates such quantities as available
bandwidth, round-trip time, server contention, the activation-transfer size
of a split, edge and server compute latency, time-to-first-token,
time-per-output-token, and resource pressure.}

\prop{%
Separating the two terms is intentional. Query-conditioned utility is a
stable property of the query--model match and is learned once; system cost is
a volatile property of the live network and hardware and is \emph{measured},
not learned. Folding $s$ into $\theta$ would require re-training whenever
network conditions change and would undermine edge deployment; we therefore
keep $C(a\mid s)$ outside the TLM weights and combine it with the learned
utility only at decision time. The scalarisation of $C(a\mid s)$ (for
example, a normalised weighted sum, or a latency penalty with the remaining
metrics reported separately, or a constrained formulation in which $\lambda$
plays the role of a budget multiplier) is treated as a measurement-and-policy
choice and does not alter the learned component $\Uhat_{\theta}$.}
% [S:08] scalarisation options S1-S4; S3/S4 would change the rule's form -> [OPEN].

% ------------------------------------------------------------------
\subsection{Training labels and utility prediction}
\label{subsec:labels}

\prop{%
The hard supervision signal is the \emph{observed} per-action outcome
$y(a,q)\in[0,1]$: for each query $q$ and action $a$ we record whether the
path produced a satisfactory answer, measured by a task-appropriate
correctness or quality score. These observed outcomes are ground truth, not
the output of any auxiliary model. Because a $\textsc{Split}(d)$ action runs
the \emph{same} model as its corresponding full execution, its quality label
equals that of the full model and only its cost differs; the split actions
therefore add no additional quality-labelling cost and are distinguished from
\textsc{Strong} purely by the measured cost term in Eq.~(\ref{eq:rule}). The
TLM is trained so that each head $\Uhat_{\theta}(a\mid q)$ is a
\emph{calibrated} estimate of $\Pr\!\bigl(y(a,q)=\text{success}\mid q\bigr)$;
calibration is enforced on held-out data and is a property we evaluate rather
than assume.}
% [S:07] y(a,q) = observed outcome (= oracle ground truth, not a teacher);
% [S:07/08] y(SPLIT(d),q) = y(full,q); target = calibrated P(success) (target E).

% ------------------------------------------------------------------
\subsection{Knowledge-distillation training}
\label{subsec:kd}

\prop{%
The TLM (the \emph{student}) is trained with a hybrid objective that combines
the hard observed outcomes with soft targets from a larger \emph{teacher}
router. The teacher is a larger encoder of the same model family, fit to the
same observed outcomes and used \emph{only during training}; it is discarded
before deployment, so it adds nothing to the inference footprint. We stress
that the observed-outcome signal $y(a,q)$ is not a teacher: it is ground
truth, and regressing the student onto it alone is ordinary supervised
learning. Distillation is the addition of the teacher's soft per-action
utilities $\Uhat_{T}(a\mid q)$, whose value is that they can generalise
beyond the labelled outcomes.}

\prop{%
Let $\alpha\in[0,1]$ weight the hard loss against the soft (distillation)
loss. The training objective is}
\begin{equation}
\label{eq:loss}
\mathcal{L} \;=\;
\alpha\, \mathcal{L}_{\text{task}}\bigl(\Uhat_{\theta},\,y\bigr)
\;+\;(1-\alpha)\,\mathcal{L}_{\text{KD}}\bigl(\Uhat_{\theta},\,\Uhat_{T}\bigr)
\;+\;\gamma\, \mathcal{L}_{\text{repr}}(\theta,T),
\end{equation}
\prop{%
where $\mathcal{L}_{\text{task}}$ is a per-action loss against the hard
targets (e.g.\ binary cross-entropy or smooth-$L_1$),
$\mathcal{L}_{\text{KD}}$ is a per-action regression of the student utilities
onto the teacher utilities, and the optional term
$\mathcal{L}_{\text{repr}}$ (weight $\gamma$, used only because student and
teacher share a family) matches an intermediate representation. Because the
primary distillation target is a per-action utility \emph{regressed}
directly, no softmax temperature is required; a temperature-scaled softmax
variant is available only as an ablation.}

\prop{%
The coefficient $\alpha$ is the central design variable, not a hidden
hyper-parameter: at $\alpha=1$ the soft term vanishes and the objective
reduces to the \emph{non-distilled} student trained on observed outcomes
alone. Knowledge distillation is therefore claimed to add value only if some
$\alpha<1$ improves held-out routing quality or calibration at a matched
student size; the non-distilled $\alpha=1$ model is the control against which
this is to be tested. We do not assert such an improvement here; it is an
empirical question, and the method remains a coherent, non-distilled
micro-encoder router if the improvement does not materialise.}
% [S:03/04] KD-3 utility regression trained as KD-5 hybrid; alpha=1 = control (E1).

\begin{algorithm}[t]
\caption{\prop{Training the distilled micro-encoder router (training-only)}}
\label{alg:train}
\begin{algorithmic}[1]
\STATE \textbf{Input:} queries $\{q_i\}$; actions $\Acal$; observed outcomes
       $y(a,q_i)$; hybrid weight $\alpha$; optional repr.\ weight $\gamma$
\STATE Fit teacher router $T$ (larger same-family encoder) on $y$; freeze $T$
\STATE Precompute soft targets $\Uhat_{T}(a\mid q_i)$ for all $i,a$
\FOR{each minibatch $B$}
  \STATE $h\gets\mathrm{Pool}(\mathrm{Encoder}_{\theta}(q))$;\quad
         $\Uhat_{\theta}(a\mid q)\gets\sigma(w_a^{\top}h+b_a)\;\;\forall a$
  \STATE $\mathcal{L}_{\text{task}}\gets$ per-action loss
         $\bigl(\Uhat_{\theta}(a\mid q),\,y(a,q)\bigr)$ \hfill(hard)
  \STATE $\mathcal{L}_{\text{KD}}\gets$ per-action regression
         $\bigl(\Uhat_{\theta}(a\mid q),\,\Uhat_{T}(a\mid q)\bigr)$ \hfill(soft)
  \STATE $\mathcal{L}\gets \alpha\,\mathcal{L}_{\text{task}}
         +(1-\alpha)\,\mathcal{L}_{\text{KD}}
         +\gamma\,\mathcal{L}_{\text{repr}}$
  \STATE $\theta\gets\theta-\eta\,\nabla_{\theta}\mathcal{L}$
\ENDFOR
\STATE Calibrate heads on held-out data (temperature / isotonic)
\STATE \textbf{Output:} student weights $\theta$ \emph{only}; discard teacher
\end{algorithmic}
\end{algorithm}

% ------------------------------------------------------------------
\subsection{Inference routing and Split-LLM integration}
\label{subsec:inference}

\prop{%
At deployment only the student $g_{\theta}$ runs, on the edge device. For a
query $q$ it computes all $N$ utilities in one forward pass; the runtime
monitor supplies the live cost $C(a\mid s)$ for each action; and the router
applies Eq.~(\ref{eq:rule}). If \textsc{Edge} is selected the small model
answers locally; if \textsc{Strong} is selected the query is served by the
remote model; if $\textsc{Split}(d)$ is selected the edge runs the first $d$
layers, transmits the activation tensor to the server, and the server runs
the remaining layers and returns the result.}

\prop{%
The router is \emph{decoupled} from the internals of the Split-LLM
implementation: it does not read or transform activations. It consumes only
(i)~a measured scalar activation-transfer size for each enabled split depth,
cached once per model, and (ii)~the live bandwidth and round-trip time, from
which the communication component of $C(\textsc{Split}(d)\mid s)$ is formed.
The split execution itself (partitioning, activation serialisation, transfer,
and server-side continuation) is provided by the existing Split-LLM runtime;
the router treats it as a measured execution path, not as a remote API call.}
% [S:09] activation-byte -> c_comm feed; Henry-code interface details are [OPEN].

\begin{algorithm}[t]
\caption{\prop{Edge inference with the distilled micro-encoder router}}
\label{alg:infer}
\begin{algorithmic}[1]
\STATE \textbf{Input:} query $q$; live system state $s$; trade-off $\lambda$
\STATE $h\gets\mathrm{Pool}(\mathrm{Encoder}_{\theta}(q))$
       \hfill(single forward pass, quantised runtime)
\STATE $\Uhat_a\gets\sigma(w_a^{\top}h+b_a)$ for each $a\in\Acal$
       \hfill(learned, calibrated)
\STATE $C_a\gets C(a\mid s)$ for each $a\in\Acal$
       \hfill(measured now, not in $\theta$)
\STATE $a^{\star}\gets\argmax_{a\in\Acal}\;[\,\Uhat_a-\lambda\,C_a\,]$
\STATE \textbf{dispatch} $a^{\star}$:
       \textsc{Edge} (local) \textbar{}
       $\textsc{Split}(d)$ (local $[0,d)$ $\to$ transfer $\to$ server $[d,L)$)
       \textbar{} \textsc{Strong} (remote)
\STATE \textbf{return} the generated response
\end{algorithmic}
\end{algorithm}

% ------------------------------------------------------------------
\subsection{Design assumptions and limitations}
\label{subsec:limits}

\prop{%
We record the assumptions on which the design rests, to be resolved
empirically rather than asserted. First, distillation is justified only if a
larger same-family teacher generalises beyond the observed outcomes; if the
teacher's soft targets coincide with the hard labels, distillation reduces to
supervised regression and the non-distilled control should be used. Second,
the number of actions $N$ and the set of enabled split depths are fixed by
the model pool and by the available Split-LLM interface, and are left open
here. Third, per-action calibration is a property we evaluate, not one the
sigmoid heads provide automatically. Fourth, Eq.~(\ref{eq:rule}) assumes a
scalar cost $C(a\mid s)$; a constrained or multi-objective treatment of the
cost would change the form of the decision rule. Fifth, all statements about
the lightness of a single forward pass are architectural; no latency has been
measured on the target edge hardware in this draft.}

% =====================================================================
% =========================  END MERGE BLOCK  =========================
% =====================================================================

\vfill
\noindent\rule{\linewidth}{0.4pt}\\
{\footnotesize\itshape
Draft provenance: this section is typeset from methods-v2 deliverable
\texttt{06-proposed-method.md} (locked f\_theta, Eqs.~3$'$/4, training and
inference algorithms, architecture) and its antecedents \texttt{00--05},
\texttt{07--09}. No experimental results are claimed. The architecture figure
referenced by the paper is specified separately in
\texttt{15-architecture-figure-spec.md} / \texttt{architecture-figure.tex}.}

\end{document}
